\scanpage{11}
\mspage{5}
même, d'après Mike Artin, prendre pour $U_i$ des schémas ``élémentaires'' sur $K_0$, se dévissant en fibrations successives de courbes. \deleted{Ceci dit, le \unsure{pt}} \inserted{(Notons que $\operatorname{Spec}K=\eta$ est le pt générique} commun des $U_i$, qui sont géom$^t$ conn. \inserted{(et $\unsure{\simeq|=}$ géom. intègres)} sur \deleted{$K_0$} $k$ (clôture alg. de $\mathbb Q$ dans $K$).

Le choix de $\bar K$ définit un pt géom. $\bar\eta$ sur les $U_i$, d'où des groupes $\pi_1(U_i,\bar\eta)=E_i$, et on a, en vertu des résultats bien connus
\[
\operatorname{Spec}K=\varprojlim U_i
\]
\[
(10)\qquad E_{\bar K/K}=\pi_1(\eta,\bar\eta)\xrightarrow{\sim}\varprojlim_i\deleted{\unsure{\pi}}E_i
\]
\margin{$(\Gamma_i=\pi_1(U_i,\bar\eta))$}
D'autre part, si on pose
\[
(11)\qquad\bar U_i=U_i\otimes_k\overline{\mathbb Q}
\]
on a \inserted{pour tt $i$} une suite exacte d'homotopie

\noindent (12)\quad\diagram{\begin{tikzcd}[ampersand replacement=\&,column sep=small,row sep=small]
1\arrow[r]\&\underbrace{\pi_1(\bar U_i,\bar\eta)}_{=\pi_i}\arrow[r]\&\pi_1(U_i,\bar\eta)\arrow[r]\arrow[d,equal]\&\pi_1(k,\bar\eta)\arrow[r]\arrow[d,equal]\&1\\
{}\&{}\&E_i\&\Gamma_{\bar k/k}=\Gamma_{\overline{\mathbb Q}/k}=\Gamma'_{\deleted{\unsure{\overline{\mathbb Q}/k}}}\&{}
\end{tikzcd}}

\margin{\diagram{\begin{tikzcd}[ampersand replacement=\&,column sep=small,row sep=small]
{}\&\operatorname{Spec}(K)=\eta\arrow[r,dash]\arrow[d]\&\eta_{\overline{\mathbb Q}}\arrow[r,dash]\arrow[d,dash]\&\bar\eta=\operatorname{Spec}(\bar K)\\
{}\&U_i\arrow[r,dash]\arrow[d,dash]\&U_{i\overline{\mathbb Q}}=\bar U_i\arrow[d,dash]\&{}\\
\mathbb Q\arrow[r,dash]\&k\arrow[r,dash]\&\overline{\mathbb Q}\&{}
\end{tikzcd}}}
qui forment un syst. projectif de suites exactes, ou d'extensions ayant toutes même quotient $\Gamma'$, et dont les noyaux
\[
\pi_i=\pi_1(\bar U_i,\bar\eta)
\]
\margin{10}
